\documentclass[11pt]{article}
\usepackage[margin=23mm]{geometry}
\usepackage[T1]{fontenc}
\usepackage{lmodern}
\usepackage{amsmath,amssymb,graphicx,booktabs,longtable,array,xcolor}
\usepackage[colorlinks=true,linkcolor=blue,urlcolor=blue,citecolor=blue]{hyperref}
\usepackage{microtype}
\setlength{\parindent}{0pt}
\setlength{\parskip}{5pt}
\newcommand{\pending}[1]{\par\noindent\fbox{\parbox{0.95\linewidth}{\textbf{Not yet available.} #1}}\par}
\newcommand{\optionalinput}[2]{\IfFileExists{#1}{\input{#1}}{\pending{#2}}}
\newcommand{\resultfigure}[3]{%
  \begin{figure}[p]\centering
  \IfFileExists{#1}{\includegraphics[width=\linewidth,height=0.71\textheight,keepaspectratio]{#1}}%
  {\fbox{\parbox[c][45mm][c]{0.91\linewidth}{\centering Numerical figure is pending.\\Expected independently generated artifact: \texttt{\detokenize{#1}}.}}}
  \caption{#2}\label{#3}\end{figure}}
\title{Alluvial ridge construction and channel--floodplain connectivity:\\an independent numerical reconstruction}
\author{Independent implementation of the written model specification of Nicholas et al. (2018)}
\date{\today}
\begin{document}
\maketitle

\begin{abstract}
The balance between channel bed aggradation and overbank deposition controls
channel depth and the pathways by which water reaches a floodplain. Nicholas
et al.\ \cite{nicholas2018} used this balance to relate alluvial ridge relief
to channel--floodplain connectivity over 580 repeated floods. This report
documents an independent reconstruction from the paper and its written
supplement. It records the physical specification, numerical assumptions,
verification evidence, and the status of all 13 manuscript panels. Numerical
panels must be generated from saved output of this implementation; published
curves and raster values are not model or plotting inputs. The current
completion status and remaining limitations are reported explicitly below.
\end{abstract}

\section{Scope and current evidence}
The calculation tests whether a coupled one-dimensional channel and
two-dimensional floodplain model yields the spatial and temporal behavior
reported by Nicholas et al.\ \cite{nicholas2018}. Their calculations describe
the period preceding an avulsion. Neither their prescribed channel migration
model nor the reconstruction should be interpreted as a prediction of an
avulsion occurrence or recurrence interval.

The permitted scientific inputs are published descriptions and configuration
values. The authors' original implementation, simulation files, plotting
scripts, and source data are excluded. Published figures were inspected for
comparison, without digitization. Published values in the supplement's
results columns are not substitutes for computed values. No published
figure is embedded in this report. There are no supplementary figures in
GSA Data Repository 2018219; DR1 is a written model description and DR2 is
a parameter and results table \cite{supplement}.

\optionalinput{report/generated/status.tex}{The panel ledger has not been rendered. Run \texttt{report/build.sh}.}

\section{Physical specification}
We model a floodplain that is 48 km wide and 72 km long. The documented
grid spacing is 250 m, giving 192 by 288 cells. The initial floodplain is
planar with downstream slope $10^{-4}$. We set its initial elevation to
$z_0(x)=10^{-4}(72\,000\,\mathrm{m}-x)$, where $x$ increases downstream
from the inlet. This vertical datum is an implementation convention.
The initially straight channel has
width $B=500$ m and depth 8 m. A flood lasts $T_f=20$ days and has minimum
and peak discharges of 2,500 and 15,000 m$^3$ s$^{-1}$. Each configuration
requires 580 floods with the same hydrograph. These values are specified
by DR1 \cite{supplement}; the hydrograph shape is not specified there.

The channel bed rises by a prescribed increment $A$ after each flood.
At any given flood it has constant slope in channel arc length $s$.
For a channel of length $L_c$, the slope is $7.2\,\mathrm{m}/L_c$, so the
inlet--outlet bed elevation drop remains 7.2 m as sinuosity changes.
For flood index $n\geq1$, the bed during the hydrograph is
\begin{equation}
 z_c(s,n)=z_{c,\mathrm{in}}^0+(n-1)A
          -7.2\,\mathrm{m}\frac{s}{L_c(n)},
 \label{eq:bed}
\end{equation}
where $z_{c,\mathrm{in}}^0=-0.8$ m and $L_c(n)$ is the channel length
during flood $n$. Saved rasters, channel coordinates, and end-flood
diagnostics precede that flood's migration, cutoff, and bed increment.
Thus a raster labeled flood $n$ contains $(n-1)A$ of imposed bed rise.
The restart checkpoint follows migration and the $n$th increment.
Equation~\eqref{eq:bed} imposes bed elevation independently of suspended
sediment exchange; it does not predict bedload transport.

The inlet concentration follows $S=CQ^{1.5}$. The supplement does not
state units for $S$ or $C$. Any conversion to SI concentration must be
identified as an assumption. The default grain fractions are 5\% sand,
75\% silt, and 20\% clay; the fine mixture is 5\% sand, 20\% silt, and
75\% clay. The bank erodibility multiplier $E$ is usually 1, with selected
values of 0.5 and 1.5. Appendix~\ref{sec:runs} lists all 31 configurations.

\subsection{Conservation laws and coupling}
The written specification requires depth-averaged shallow-water flow and
three suspended sediment fractions. In a floodplain cell, let $h$ be water
depth, $\mathbf{u}=(u,v)$ velocity, $z$ bed elevation, and $c_k$ the suspended
mass concentration for grain class $k$. Away from external boundaries,
the governing balances have the form
\begin{align}
 \partial_t h+\nabla\cdot(h\mathbf{u}) &= 0, \\
 \partial_t(h\mathbf{u})+
 \nabla\cdot\left(h\mathbf{u}\otimes\mathbf{u}
                   +\tfrac12 gh^2\mathbf{I}\right)
 &= -gh\nabla z-\mathbf{f}_b+\nabla\cdot\mathbf{T}, \label{eq:momentum}\\
 \partial_t(hc_k)+\nabla\cdot(h\mathbf{u}c_k)
 &= \nabla\cdot(hD_k\nabla c_k)-d_k . \label{eq:sediment}
\end{align}
Here $\mathbf{f}_b$ is bed resistance per unit water density,
$\mathbf{T}$ represents any chosen horizontal momentum diffusion,
$D_k$ is sediment diffusivity, and $d_k$ is deposition per unit bed area
and time. The written sources determine the required transport structure;
the final constitutive choices and any omitted terms must be stated in the
implementation record below. If concentration is in kg m$^{-3}$, $d_k$ is
in kg m$^{-2}$ s$^{-1}$. Deposit thickness is related to deposited mass by
\begin{equation}
 \partial_t z=\frac{\sum_k d_k}{\rho_s(1-p)},
 \label{eq:deposit}
\end{equation}
where $\rho_s$ is sediment density and $p$ is deposit porosity. Equation
\eqref{eq:deposit} applies to the modeled depositional bed changes, not to
the separately prescribed channel aggradation.

The channel uses section-integrated water and momentum equations and
advection--diffusion of three suspended fractions, with lateral exchange
and boundary supply. Net exchange between channel suspension and its
prescribed bed is set to zero. Channel aggradation and geometry adjustments
are independent signed bed forcing, recorded separately in the solid-mass
budget. DR1 prescribes the bed and excludes bedload-based updates, but does
not establish zero net exchange with suspension \cite{supplement}.
This is therefore a physical assumption, not a recovered closure.
Channel variables are interpolated to floodplain-grid cells within
$B/2$ of the centerline, with velocity rotated into the local channel tangent.
Interface face fluxes update the floodplain and the corresponding channel
section with equal and opposite water and sediment mass transfers. The
overlaid channel cells must not duplicate the water or sediment storage
already represented by the one-dimensional channel.

At the end of each flood, the written specification calls for
Howard--Knutson migration with neck cutoffs \cite{howard1984}. Newly occupied
channel cells are lowered to the interpolated channel bed. Their eroded
sediment is returned at the upstream concentration boundary. The timing,
grain allocation, and mass accounting for this return are implementation
choices where DR1 is incomplete. The topography of abandoned channel cells,
cutoff handling, and remapping of water and sediment also require explicit
choices. A migration displacement is not itself a shallow-water avulsion.

\subsection{Implementation and unresolved assumptions}
\optionalinput{report/generated/implementation.tex}{The full coupled implementation record is pending. The initial numerical kernel uses a first-order HLL flux with hydrostatic reconstruction and a fixed-depth, well-mixed settling substep. These choices alone do not implement the complete coupled model. See \texttt{docs/specification.md} for the current assumption ledger.}

The horizontal Boussinesq turbulent stress is implemented, but the
secondary-flow and associated dispersive terms in the earlier methods
paper are omitted \cite{nicholas2013}. These terms describe effects of
unresolved vertical velocity structure; horizontal eddy viscosity does
not replace them. Bed composition is also not tracked as deposits
accumulate and are reworked. Excavated bank mass is assigned the selected
inlet grain fractions. This preserves total solid mass while losing the
grain identity of earlier deposits. Together with omitted vertical
floodplain entrainment and zero net channel suspended-bed exchange, these
choices limit correspondence with the inherited process equations.

DR1 specifies first-order spatial accuracy and an explicit finite-volume
approximate Riemann solver. The independent kernel uses hydrostatic
reconstruction \cite{audusse2004}. At an interface, $z_* = \max(z_L,z_R)$ and
$h_* = \max(0,h+z-z_*)$. Velocities and concentrations are retained in the
reconstructed wet states. Side-specific pressure corrections balance a
stationary free surface over bed steps. The HLL water flux and the sediment flux formed from upwind concentration
are shared between the two sides. These are independently selected numerical
details, not a claim to match the original implementation.

\section{Diagnostics and verification}
The published super-elevation ratio is
\begin{equation}
 \beta_n=\frac{z_{\mathrm{ridge},n}-z_{\mathrm{distal},n}}
                  {H_{\mathrm{channel},n}}.
\end{equation}
The reported late-time statistic is the arithmetic mean of the 20
per-flood ratios $\beta_n$ at floods 561--580. Ridge relief is the maximum
transverse value of longitudinal-mean relief over the central 80\% of the
reach, relative to the outermost 10\% lateral bands. Channel depth is the
mean bank-face elevation minus adjacent channel-bed elevation. These
spatial averages and the mean-of-ratios convention are implementation
choices not resolved by the supplement. A constant vertical datum cancels
from the numerator, but changes in the selected ridge or distal region do
not.

For each flood, the connectivity measure is sampled at $t=0.6T_f=12$ days:
\begin{equation}
 \alpha_n=\frac{\overline{h}_{\mathrm{fp},n}}
                         {f_{\mathrm{wet},n}},\qquad
 \alpha_{95}=P_{95}\{\alpha_{381},\ldots,\alpha_{580}\}.
 \label{eq:alpha}
\end{equation}
Here $f_{\mathrm{wet}}$ is the inundated fraction of the channel--floodplain
interface. Unlike $\beta$, $\alpha$ has units of meters. The depth average
includes every floodplain cell, including dry cells. The interface consists
of internal Cartesian faces between active channel and floodplain cells;
all faces have equal length at the uniform grid spacing. A face is
inundated when floodplain depth and channel overtopping depth both exceed
1 mm. Zero inundated fraction gives undefined $\alpha$. The percentile
uses linear interpolation and requires all 200 values to be finite;
missing or undefined values are neither replaced by zero nor discarded.

At day 12 the output retains $\alpha$, mean floodplain depth, and
inundated interface fraction for every flood. For the requested map and
profile floods it also saves rasters and channel coordinates with a
\texttt{peak\_} prefix, at cumulative time $(n-0.4)T_f$. These files
support assessment of inundated area and breach location at the sampling
time, which deposit footprint alone cannot supply. The earlier interrupted
rebuild did not save these fields. Their availability and spatial
interpretation remain to be assessed from the new calculations. Channel
geometry is unchanged within each hydrograph, so the end-flood outline
identifies the full day-12 interface, without identifying its inundated
subset.

Verification must distinguish discrete conservation from agreement with
the paper. Required checks include a lake at rest, positive wetting and
drying, coupled water and grain-mass exchange, analytical transport and
settling, zero forcing limits, and cutoff/remapping budgets. External
boundary fluxes, bed changes, and prescribed channel aggradation must be
reported separately. Representative coupled cases require grid and time
step sensitivity. Kernel tests alone cannot establish these properties for
the full coupled model.

\optionalinput{report/generated/validation.tex}{Full coupled verification and grid/time sensitivity are pending. The kernel check output, if present, is transcribed below without extending its scope.}
\optionalinput{report/generated/kernel.tex}{No kernel verification log was available at report build time.}

Interpretation of production conservation and sensitivity remains pending.
Finite residuals alone do not establish negligible budget drift. The final
assessment requires maximum and terminal residuals with stated
normalization, together with the numerical field differences and diagnostic
changes from grid, time-step, and physical-assumption cases. Per-flood
budget summaries precede migration and aggradation. A separate
\texttt{remap\_budget.csv} records local water and suspended-class
transfer defects, the combined solid defect after signed bed forcing,
and cumulative residuals after each remap. Both cumulative ledgers reset
their reference inventory at process restart; they do not recover
accumulated full-history drift across restarts. The sensitivity cases
cover the first flood and retain 250 m channel spacing,
so they cannot establish 580-flood migration or grid independence.

\section{Independent figures and comparison targets}
Figures~\ref{fig:deposits}--\ref{fig:diagnostics} reserve all manuscript panels.
A placeholder is a missing result. A generated panel must retain the
configuration, sampling time, and plotted-data provenance in
Appendix~\ref{sec:panels}. Status is assessed separately from whether a
plotting script produced a file: a valid numerical panel may still disagree
with the publication.

\subsection{Flood deposits: manuscript panels 1A--F}
The published comparison is between deposits from individual floods, not
accumulated bed elevation. Panels 1A--D use $A=0$, $C=0.003$, and the default
load at floods 61, 251, 501, and 562. Panels 1E and 1F use flood 562, with
$(A,C)=(0.036,0.003)$ and $(0,0.0003)$, respectively. The published color
scale spans 0--0.5 m, with deposits below $10^{-6}$ m distinguished from
the active channel. This scale is a comparison reference, not a clipping
rule for the computed mass budget.

The paper reports early splay deposition near migrating bend apices,
progressive widening of the channel belt, and later deposition along
abandoned channels. It reports reduced floodplain inundation as channel
depth increases, with event-to-event fluctuations under identical imposed
floods. The high-aggradation and lower-load cases retain more widespread
inundation. Independent comparison should measure deposit area above
$10^{-6}$ m, deposit volume, transverse spread, channel sinuosity, and
abandoned-channel coverage where saved output supports these diagnostics.
No pixel-wise match is expected without the undocumented initial migration
perturbation and dimensional migration parameters.

\resultfigure{figures/figure1.pdf}{Independent model output for manuscript
panels 1A--F: deposit thickness during each listed flood, in meters, on the
48 km by 72 km domain. A--D use $A=0$, $C=0.003$, $E=1$, and the default
grain mixture, at floods 61, 251, 501, and 562. E uses $A=0.036$ and
$C=0.003$ at flood 562; F uses $A=0$ and $C=0.0003$ at flood 562.
Flow is toward increasing downstream distance. Black marks active channel
cells and green marks deposits below $10^{-6}$ m. The color scale extends
to 0.5 m; larger values saturate visually but remain in the saved data and
deposit-volume calculation. The solid magenta line in E follows the full
active channel--floodplain interface used for the connectivity calculation,
not its inundated subset. Snapshots precede each flood's migration and
aggradation. These panels test whether
levee breaches and abandoned channels focus deposition as the ridge grows.
The panel ledger gives the status and commands for each requested result.}{fig:deposits}

\subsection{Transverse elevation: manuscript panels 2A--D}
The cross-section uses the downstream one of the two saved grid columns
equidistant from the domain midpoint, at 36.125 km for the 250 m grid.
Its initial floodplain elevation is 3.5875 m under the model datum.
Each panel requests elevations at the end of floods 50, 225, 400,
and 580, before migration and aggradation. The published panels show a
growing ridge with a concave cross-sectional envelope, interrupted by
active and abandoned channels.
The approximate late-time ridge elevations visible in the published
panels are 40--45 m (2A), 25--30 m (2B), 15--17 m (2C), and 30--33 m (2D).
These are coarse visual ranges, not digitized observations. The initial
mid-domain floodplain is plotted near 9 m in the publication; DR1 leaves
the absolute vertical datum unspecified. The independent plots apply no
datum shift. Quantitative comparisons subtract each initial plane and
retain the uncertainty of the published visual range. Channel minima,
ridge width, and the sign and scale of relief growth provide additional
comparisons.

\resultfigure{figures/figure2.pdf}{Independent model output for manuscript
panels 2A--D: floodplain elevation versus transverse distance at the
central grid column, 36.125 km downstream, at the end of floods 50, 225,
400, and 580, before migration and aggradation. Black solid, gray dotted,
orange dashed, and red solid lines identify those floods, respectively.
Cyan marks saved active channel-bed cells at flood 580. All cases use
$E=1$ and the default grain mixture. A: $A=0.036$, $C=0.01$;
B: $A=0.036$, $C=0.0003$;
C: $A=0$, $C=0.0003$; D: $A=0.018$, $C=0.003$.
Aggradation $A$ is in m per flood. Distance and elevation are in km and m,
respectively. The initial section is at 3.5875 m; no datum shift is applied.
The comparison tests ridge relief growth and channel deepening under different balances
of prescribed bed aggradation and computed overbank deposition.}{fig:sections}

\subsection{Super-elevation and connectivity: manuscript panels 3A--C}
Panel 3A requires one late-time $\beta$ value for each of the 31
configurations. The paper reports that $\beta$ rises with $A$ and, over
most shared $A$ values, decreases with increasing $C$. A finer load reduces
ridge super-elevation. The prose reports approximately a 15\% increase in
$\beta$ for lower erodibility and a 6.5\% decrease for higher erodibility;
these are comparisons to test, not coefficients to impose. Only one of
the fine-load simulations has published $\beta>1$. The $\beta=1$ line is
an interpreted topographic threshold, not a modeled avulsion criterion.

Panel 3B compares $\alpha_n$ for $C=0.0003$, 0.003, and 0.01 at
$A=0.018$ m per flood. The paper reports stronger temporal variability
and higher peaks as sediment load increases. The published axis spans
approximately 0.1--10 m; exact event timing is not identifiable from the
written specification. Panel 3C compares $\alpha_{95}$ across all runs.
The main reported pattern is large values for high suspended load and
low aggradation, when conveyance occurs through localized breaches.
Numerical comparison should report the computed ranges, trend signs,
late-time variability, and any failure to preserve this ordering. A high
$\alpha$ by itself does not demonstrate avulsion.

\resultfigure{figures/figure3.pdf}{Independent model output for manuscript
panels 3A--C. A: arithmetic mean of the per-flood super-elevation ratios
$\beta_n$ over floods 561--580, against imposed bed aggradation $A$, for
all 31 configurations.
B: connectivity $\alpha$ at day 12 of each flood for $A=0.018$ m per flood
and $C=0.0003$, 0.003, and 0.01, with default load and $E=1$.
C: linearly interpolated 95th percentile of $\alpha$ over all 200 floods
381--580 against $A$.
$\beta$ is dimensionless; $\alpha$ and $\alpha_{95}$ are in meters.
Purple, blue, orange, and black identify $C=0.00003$, 0.0003, 0.003,
and 0.01, respectively. Squares denote default load with $E=1$,
diamonds denote fine load, and plus and cross symbols denote $E=0.5$
and 1.5. Red rings mark altered or unverified settings or incomplete
histories. The dashed line at $\beta=1$ is a published interpretive
reference. Missing values create gaps; an incomplete or nonfinite
statistical window yields no point. Alpha axes are logarithmic when all
available values are positive. The panel ledger records current coverage
and provenance.}{fig:diagnostics}

\clearpage
\optionalinput{report/generated/results.tex}{Quantitative comparison from completed coupled simulations is pending. The comparison targets above describe the publication and do not constitute independently reproduced results.}

\section{Limits of the reconstruction}
Assessment of the generated morphology remains pending. The automatic
comparisons report scalar spatial scales, trend signs, and variability,
but do not establish splay formation, abandoned-channel infilling, ridge
concavity, or breach localization. Those claims require direct assessment
of the independent maps and profiles and an authored record for each
panel. Deposit footprint is not inundation extent. Per-flood alpha
components and cutoff counts support temporal comparisons. Day-12 spatial
evidence is limited to the requested snapshot floods, and the original
initial migration perturbations are undocumented.
Completion of all numerical windows therefore does not, by itself,
establish scientific reproduction.

The largest identifiable ambiguities concern the sediment concentration
units, grain diameters and settling closures, flow resistance, the
hydrograph shape, and the dimensional migration scale. These choices
change sediment delivery, bankfull capacity, and the time needed to widen
the channel belt. Their effects can therefore propagate into every panel.
The definitions of ridge height and inundated interface add uncertainty
to $\beta$ and $\alpha$ even when the underlying hydraulic solution is
unchanged. Assumption sensitivity must be separated from numerical
resolution sensitivity.

The written model fixes channel width and bed aggradation. It does not
resolve the bedload supply required to sustain the imposed bed elevation
changes. Comparisons consequently concern alluvial ridge construction and
connectivity within that idealization. Inferences about natural avulsion
likelihood remain bounded by those assumptions.

\section{Rebuild and deliverables}
\optionalinput{report/generated/rebuild.tex}{A verified clean model-to-report rebuild remains pending. To rebuild this report from the current ledger and figures, run \texttt{bash report/build.sh} from the repository root. That command does not run missing simulations and does not establish a completed reconstruction.}

\appendix
\section{Figure-by-figure record}\label{sec:panels}
The record is generated from \texttt{docs/panels.csv}. Every numerical
panel requires both a data command and a plot command. ``Not reproduced''
includes unattempted or failed computations; ``partially reproduced'' must
name the missing sampling, physics, resolution, or agreement; ``independently
reproduced'' is reserved for a documented panel whose plotted values come
from this implementation. A schematic cannot satisfy a numerical panel.
\optionalinput{report/generated/panels.tex}{Panel table has not been generated.}

\section{Configuration matrix}\label{sec:runs}
This table contains input parameters only. It is generated from
\texttt{config/runs.csv}, transcribed from the configuration columns of
DR2 \cite{supplement}. $A$ is in m per flood. Units for the published $C$
are unspecified. All runs require 580 floods at the documented resolution.
\optionalinput{report/generated/runs.tex}{Configuration table has not been generated.}

\begin{thebibliography}{9}
\bibitem{nicholas2018} Nicholas, A. P., Aalto, R. E., Sambrook Smith,
G. H., and Schwendel, A. C. (2018). Hydrodynamic controls on alluvial ridge
construction and avulsion likelihood in meandering river floodplains.
\emph{Geology}, 46, 639--642. \url{https://doi.org/10.1130/G40104.1}.
\bibitem{supplement} Nicholas, A. P., Aalto, R. E., Sambrook Smith,
G. H., and Schwendel, A. C. (2018). GSA Data Repository 2018219,
DR1: Model description; DR2: Table of model parameters and results.
\url{https://doi.org/10.1130/2018219}.
\bibitem{howard1984} Howard, A. D., and Knutson, T. R. (1984).
Sufficient conditions for river meandering: A simulation approach.
\emph{Water Resources Research}, 20, 1659--1667.
\url{https://doi.org/10.1029/WR020i011p01659}.
\bibitem{audusse2004} Audusse, E., Bouchut, F., Bristeau, M.-O.,
Klein, R., and Perthame, B. (2004). A fast and stable well-balanced scheme
with hydrostatic reconstruction for shallow water flows.
\emph{SIAM Journal on Scientific Computing}, 25, 2050--2065.
\url{https://doi.org/10.1137/S1064827503431090}.
\bibitem{nicholas2013} Nicholas, A. P., Ashworth, P. J., Sambrook Smith,
G. H., and Sandbach, S. D. (2013). Numerical simulation of bar and island
morphodynamics in anabranching megarivers.
\emph{Journal of Geophysical Research: Earth Surface}, 118, 2019--2044.
\url{https://doi.org/10.1002/jgrf.20132}.
\bibitem{ferguson2004} Ferguson, R. I., and Church, M. (2004).
A simple universal equation for grain settling velocity.
\emph{Journal of Sedimentary Research}, 74, 933--937.
\url{https://doi.org/10.1306/051204740933}.
\bibitem{schwendel2015} Schwendel, A. C., Nicholas, A. P., Aalto, R. E.,
Sambrook Smith, G. H., and Buckley, S. (2015). Interaction between meander
dynamics and floodplain heterogeneity in a large tropical sand-bed river:
the Rio Beni, Bolivian Amazon.
\emph{Earth Surface Processes and Landforms}, 40, 2026--2040.
\url{https://doi.org/10.1002/esp.3777}.
\bibitem{limaye2014} Limaye, A. B. S., and Lamb, M. P. (2014).
Numerical simulations of bedrock valley evolution by meandering rivers
with variable bank material.
\emph{Journal of Geophysical Research: Earth Surface}, 119, 927--950.
\url{https://doi.org/10.1002/2013JF002997}.
\end{thebibliography}
\end{document}
